\chapter{Applications in Biology and Ecology}
\label{ch:applications}

A survey of where these methods have actually produced results, and where the
claims outrun the evidence. The distinction matters: the successes and the
overreaches have a recognisable shape, and learning to tell them apart is
worth more than memorising either list.

\section{The pattern behind the successes}

Structure prediction is the canonical case. AlphaFold worked because protein
folding had the properties a deep learning problem needs: a large corpus of
solved structures as ground truth, a well-defined input--output mapping, a
physically constrained output space, and an unambiguous evaluation metric. The
model learned a real regularity because a real regularity was there and was
densely sampled.

\textcite{brixi2026evo2} extend the same logic to genomes --- modelling and
design across all domains of life, at scale, from DNA sequence. Again: enormous
training corpus, discrete well-defined alphabet, evaluable output.

\textcite{jiang2025aibioinformatics} survey the broader field and report the
landmark numbers: near-atomic protein structure prediction at median 0.96\,\AA{}
on CASP14, protein design success rates up to 92\%, sensitive cancer detection
at AUC $\approx$ 0.93.

Deep learning succeeds where there is a large, densely sampled, well-labelled
corpus and an unambiguous target. Ecological field data rarely has any of the
three.

\section{What the ecology literature shows}

Reviews of language models written for ecologists are rare.
\textcite{mammides2026llmecology} is the fullest so far: 129 documents on
ecology and conservation (85 peer-reviewed articles, 38 preprints, 6
conference papers), sorted by what the models were used for. The preprint share
tells you how young the evidence is. The main uses they found, with an example
of each:

\begin{description}
  \item[Extracting data from text] Plant pathogens and their hosts from disease
    reports, faster than human reviewers and with high accuracy; species
    interactions from over 80{,}000 articles; wildlife-crime hotspots for sea
    turtles from 247 Chinese court cases. Accuracy varied a great deal with
    model and task.
  \item[Querying databases in plain language] A chatbot over iDigBio, and a
    pipeline that answers questions about eBird observations.
  \item[Literature synthesis] Over two million articles screened for evidence
    on nature-based solutions; literature reviews run as repeated prompting
    with manual checks.
  \item[Monitoring public sources] NewsPanda, a fine-tuned classifier used by
    WWF to track news about thousands of conservation sites; detection of
    wildlife trafficking in online adverts; visitor reviews of protected areas.
  \item[Code] Writing, explaining, commenting and translating analysis code,
    notably from R to Python, for a field where formal programming training is
    often thin.
  \item[Images and sound] BioCLIP, trained on 10 million images covering over
    454{,}000 taxa, identifies species it was not trained on; WildCLIP finds
    rare events in Serengeti camera-trap data; NatureLM-audio describes animal
    calls from plain-language prompts.
  \item[Communication and policy] Translation and editing for non-native
    English speakers; OceanGPT, retrained on 67{,}000 ocean science documents;
    checking national biodiversity targets against the global ones; and,
    speculatively, agents that simulate stakeholders to test a policy.
\end{description}

Their list of problems is as long as the list of uses. Hallucination is worst
on uncommon topics, which in ecology means most of them. Answers carry the bias
of the training data: asked about ecological restoration, ChatGPT favoured
researchers in high-income countries and underrepresented lower-income regions
and Indigenous communities. Outputs are hard to reproduce because the models are
opaque, give different answers to the same prompt, and are withdrawn without
notice. Paid access and the skills needed may widen the gap between rich
countries and poorer, often species-rich, ones. Relying on the model for code,
synthesis and writing risks \term{deskilling} --- losing the skills needed to
check it. And there is the running cost in energy and water.

One warning in the review deserves repeating for this audience. Field studies
are already a shrinking share of the literature, and students spend less time
in the field. If language models pull attention and money further away from
collecting primary observations, which are what ecological knowledge rests on,
the authors count that as a major loss. It is the same point
\textcite{mishra2026notready} makes from the data side.

Their recommendations are the ones this book argues for: choose the model for
the task, write careful prompts, ground answers with retrieval, keep a person
checking the output, and publish prompts and workflows with the results.

\section{Where ecology has genuine wins}

\subsection{Image classification}

Camera traps, acoustic monitoring and remote sensing are the strongest cases,
and for the same reason as structure prediction: large labelled corpora and an
unambiguous target. Species identification from camera trap imagery routinely
exceeds 90\% accuracy for well-represented taxa while collapsing manual
annotation time \parencite{mishra2026notready}. Satellite imagery classification
for forest cover change and vegetation type operates at spatial and temporal
resolutions that were previously unattainable.

The caveats are the usual ones and they are not minor. Accuracy is reported on
common species; rare species --- the ones conservation actually cares about ---
are underrepresented in training data and classified worst. Models trained in
one biogeographic region transfer poorly to another. And a classifier confident
about a species outside its range is a data quality problem that propagates
silently into downstream analysis.

\subsection{Natural history collections}

\textcite{johnson2026collections} report robotic imaging combined with AI-driven
transcription of specimen labels at scale. This is a strong application because
the bottleneck was always transcription throughput, not scientific judgement,
and because the output is verifiable against the physical specimen. Roughly a
billion specimens sit in the world's collections with label data locked in
handwriting and typescript. Mobilising that is a mass-transcription problem, and
mass transcription is what these systems do.

\subsection{Text mining the literature}

Biodiversity knowledge is disproportionately locked in prose --- in taxonomic
descriptions, expedition reports, faunal lists, and the grey literature.
\textcite{castro2024unstructured} demonstrate LLM extraction from exactly this
material.

\textcite{takola2026topicmodelling} traces the methodological arc for
conservation specifically: Latent Semantic Analysis to Latent Dirichlet
Allocation to embedding-based models like BERTopic and Top2Vec to LLMs. Her
conclusion deserves attention because it runs against the upgrade narrative:
traditional probabilistic models remain valuable. They are interpretable,
cheap, reproducible and auditable. An LLM topic assignment is none of those
things. Newer is not automatically better when what you need is a method section
you can defend.

Bibliometric synthesis is a related and undersold application.
\textcite{ioannou2025onehealth} map One Health research across environmental,
animal and human health dimensions in coastal and marine contexts --- the kind
of cross-domain landscape analysis that is tedious manually and tractable
computationally.

Infrastructure is appearing for the retrieval side. Biodiversity PMC
\parencite{pasche2026biodivpmc} indexes biodiversity literature specifically.
\textcite{krug2026ipbes} examine the role of AI in IPBES assessments and pose the
question in their title --- hype or salvation --- without resolving it, which is
the right answer for assessment work, where every statement has to be traceable
to a source.

\subsection{Microbiome and multi-omics}

\textcite{yu2026microbiomes} review the move from descriptive microbiome
analysis toward predictive modelling and hypothesis generation, covering
integrative analysis of microbial genomes, community structure and functional
potential. They are explicit about what blocks the field: model
interpretability, generalisability across cohorts, and causal inference. The
first two are engineering problems. The third is not.

\subsection{Curation support}

\textcite{pombase2026curation} describe the application closest to ELIXIR's own
operations: AI-assisted literature curation for biological knowledgebases.
Their framing is the correct one --- \term{human-in-the-loop}, with generative
methods accelerating routine tasks (recommendation, entity recognition,
extraction, summarisation, ontology development, quality control) while curators
retain critical evaluation. They identify the risks precisely: scarce training
data, difficulty extracting complex relationships, and error propagation into
resources that other work depends on.

The last one is the reason this must stay human-supervised. A knowledgebase error
does not stay local; it propagates into every analysis downstream, and
attribution back to its source becomes impossible once it has been cited.

\subsection{Query generation}

\textcite{smeros2025sparqlllm} address a real friction point: writing SPARQL
against federated life science knowledge graphs is a specialist skill that
excludes most domain scientists. Their system generates queries from natural
language questions and --- unusually, and to their credit --- evaluates on
federated query capability, runtime and cost rather than accuracy alone.

This pattern generalises well. Generating a query is safe in a way that
generating an answer is not: the query is inspectable before execution, the
database provides ground truth, and a wrong query usually fails loudly. Prefer
query generation over answer generation wherever the data is in a queryable
store.

\section{Where the claims outrun the evidence}

\textcite{mishra2026notready} makes the case directly, and the argument is worth
stating in full because it is the most important dissent in this literature.
Ecological data is heterogeneous, spatially and temporally biased, sparse for
the taxa that matter most, and generated by processes that are causal rather
than merely correlational. A model that predicts species occurrence from
environmental covariates learns the sampling effort as readily as the ecology.
Deploying such models for conservation decisions transfers the bias into policy
while adding a veneer of quantitative authority.

The specific hazards:

\begin{itemize}
  \item \textbf{Sampling bias passed off as signal.} Occurrence data records
    where people looked. A model trained on it predicts where people looked.
  \item \textbf{Absent causal structure.} Correlation between a covariate and a
    distribution is not a mechanism, and extrapolating to novel climate
    conditions requires the mechanism.
  \item \textbf{Uncertainty that does not propagate.} Deep learning outputs are
    typically point predictions. Ecological inference needs intervals, and
    downstream decisions need them honestly propagated.
  \item \textbf{Unfalsifiable models.} An uninterpretable model that fits well
    cannot be distinguished from one that fits well for the wrong reason.
\end{itemize}

\begin{keypoint}
Before applying a model to an ecological question, ask what would have to be true
of the data for the answer to be valid --- and whether it is. The question is not
whether the method is powerful. It is whether the data can support the inference.
\end{keypoint}

\section{Agents in the life sciences}

A distinct category: not models that classify, but agents that run analyses.
Biomni \parencite{huang2025biomni} is a general-purpose biomedical agent;
PantheonOS \parencite{xu2026pantheonos} a multi-agent framework for genomics
discovery; OmicClaw \parencite{zeng2026omicclaw} targets executable and
reproducible natural-language multi-omics analysis. Domain-specific skill
libraries are emerging alongside them.

The evaluation literature is the useful part.
\textcite{guo2026promptbio} benchmarked these systems on 244 expert-curated
bioinformatics tasks and found comparable performance between leading agents and
sharp decline with difficulty. The honest reading: useful for scoped,
well-specified analytical tasks; unreliable for anything requiring sustained
judgement about whether the analysis is the right one.

The most striking result so far comes from inside a model company.
\textcite{yoon2026art}, at Anthropic, set Claude Code agents loose on reverse
transcriptases --- enzymes that copy RNA into DNA --- across 1.9 billion protein
clusters. Worker agents planned and ran each task, supervisor agents reviewed
them and opened new ones, and all of it went into a shared record: 119 tasks
and 949 agent sessions over 21.5 hours, with no human intervention. Of 17
candidate partner genes, 14 turned out to be annotation artefacts or parts of
known systems. One worker, reading the raw DNA, noticed a run of
repeats beside a phage enzyme. That became a new family, array-associated
reverse transcriptases, whose repeat arrays resemble CRISPR arrays.

Two caveats. The paper comes from the model's makers, using their own frontier
model. And, as the authors say, the finding is a prediction from sequence
alone: they have not shown that the enzyme is active or what the system
does for the phage. The agents noticed an anomaly, which is where discovery
starts; confirming it is still bench work.

Nor is the randomness solved. Two runs of the same agent on the same input can
give different answers. If the workflow has to be reproducible, that is a
serious problem rather than a rough edge --- see \Cref{ch:reproducibility}.

\section{Domain assistants}

UNEP's EnvironmentGPT \parencite{unep2026environmentgpt} is instructive as a
design pattern: a RAG system restricted to a curated library of UNEP and partner
publications, rather than a general model over uncontrolled web text. The
tradeoff is explicit --- narrower coverage in exchange for citable answers and
reduced hallucination.

\textcite{sigillo2026embllibrarian} generalise this into infrastructure: a
life-sciences knowledge layer that AI agents query, rather than each agent
scraping and indexing independently. This is the right architecture. The
alternative, in which every group builds its own RAG pipeline over the same
literature, wastes effort and produces inconsistent answers to identical
questions.

For an assistant of this kind, a curated set of documents with citations beats a
larger model without them. The value is in the documents and in being able to
check where an answer came from, not in the model.
